\section{Introduction}\label{sec:introduction}

Haldane predicted that antiferromagnetic Heisenberg chains have different
low-energy behavior for integer and half-odd-integer spins. The prediction
was first formulated in 1981 and published in 1983
\cite{Haldane1981,Haldane1983PRL,Haldane1983PLA}. For spin one, the prediction is a
positive gap above the ground energy. We prove this prediction for the
pure bilinear model on even periodic rings, with explicit bounds uniform
in the length.

Let $S^x,S^y,S^z\in M_3(\C)$ be the self-adjoint irreducible spin-one
matrices, normalized by
\[
 [S^\alpha,S^\beta]
   =i\sum_\eta\varepsilon_{\alpha\beta\eta}S^\eta,
 \qquad \sum_\alpha(S^\alpha)^2=2I_3.
\]
For each even integer $L\ge4$, let $\cH_L=(\C^3)^{\otimes L}$ and define
\begin{equation}\label{eq:main-H}
 H_L=\sum_{j=0}^{L-1}\sum_{\alpha=x,y,z}
       S_j^\alpha S_{j+1}^\alpha,
 \qquad S_L^\alpha=S_0^\alpha.
\end{equation}
Here $S_j^\alpha$ acts on site $j$ and is the identity on the other
factors. The coupling constant is one. Write $E_0(L)$ for the least
energy and $E_1(L)$ for the next \emph{distinct} energy of $H_L$, and set
\begin{equation}\label{eq:main-gap}
 \gamma_L=E_1(L)-E_0(L),
 \qquad
 \Delta_1=\liminf_{\substack{L\to\infty\\L\text{ even}}}\gamma_L.
\end{equation}
These definitions concern the full spectrum, including the multiplicity
of the ground energy. The next distinct energy exists: the all-zero and
all-$+1$ product vectors in the $S^z$ basis have energy expectations
$0$ and $L$, respectively, so $H_L$ is not scalar.

\begin{theorem}\label{thm:main}
For the Hamiltonians \eqref{eq:main-H}, the ground state is unique for
every even $L\ge60$, and
\begin{align}
 \gamma_L&>\frac4{105}\log\frac{80}{79}
     &&\text{for every even }L\ge60,\label{eq:gap60}\\
 \gamma_L&>\frac{\log20}{784}
     &&\text{for every even }L\ge2304.\label{eq:gap2304}
\end{align}
Consequently
\[
 \Delta_1\ge\frac{\log20}{784}>0.
\]
The spectral conclusions also hold for any simultaneously similar
spin-one realization satisfying the same commutator and Casimir
normalizations.
\end{theorem}

Theorem~\ref{thm:main} resolves the spin-one Haldane gap conjecture
positively in the even-periodic finite-volume formulation
$\Delta_1>0$. It also implies $\inf_{L\ge4,\,L\text{ even}}\gamma_L>0$:
each of the finitely many remaining finite-dimensional Hamiltonians has
a positive gap above its full ground sector. This observation supplies
neither a numerical value nor a simplicity assertion at those small
lengths.
Corollary~\ref{cor:local-limit} also gives the corresponding lower bound
for local excitations in every subsequential thermodynamic limit of the
periodic ground states of the self-adjoint model.

\subsection{Historical context}

The half-odd-integer side has a rigorous antecedent in the low-energy
excitation bound of Lieb, Schultz and Mattis for the cyclic spin-one-half
chain~\cite[Appendix~B, Theorem~2]{LiebSchultzMattis1961}.
Affleck and Lieb extended this obstruction to the half-odd-integer
antiferromagnetic chains they studied~\cite{AffleckLieb1986}. The possible alternative of degenerate
ground states is part of that obstruction. Integer spin permits a unique
gapped state, but this distinction alone does not prove a gap for a
particular interaction.

The Affleck--Kennedy--Lieb--Tasaki construction established a rigorous
gapped spin-one chain with a biquadratic interaction
\cite{AKLT1987,AKLT1988}. Its bond interaction is
$h+\tfrac13h^2$, up to a scalar and normalization, where
$h=\mathbf S_j\cdot\mathbf S_{j+1}$. The Hamiltonian in
\eqref{eq:main-H} has bond interaction $h$ itself.
Yarotsky proved that a unique ground state and a uniform periodic gap
persist under sufficiently small translation-invariant finite-range
perturbations of the AKLT interaction
\cite[arXiv v1, Theorems~1 and~3]{Yarotsky2006}.
This controls a neighborhood of that interaction; it does not establish
a gapped path to the pure bilinear point.

For the pure spin-one chain, White and Huse's density-matrix
renormalization-group study
estimated a bulk gap of $0.41050(2)$ and exhibited effective spin-one-half
degrees of freedom at open ends~\cite{WhiteHuse1993}. Neutron scattering
provided experimental evidence in the nearly one-dimensional spin-one
antiferromagnet $\mathrm{CsNiCl}_3$~\cite{Buyers1986}.
These results explain both the strong
physical evidence for the prediction and the significance of boundary
conditions. The uniform lower bounds below instead follow from the
stated finite certificates and an analytic argument valid at every
larger even length.

Batista and Ortiz proposed a spin--fermion proof of the Haldane gap
\cite{BatistaOrtiz2001}. Their gap argument following Equation~(8) of the preprint version
uses a first-order perturbative identification with an XYZ chain;
that step does not supply a uniform remainder estimate or a continuation
to the original pairing coefficient. More recently, Tasaki established
a nontrivial infinite-chain topological index for the spin-one
Heisenberg model under a gap assumption for odd open chains with endpoint fields
\cite[arXiv v4, Assumption~2 and Theorem~3]{Tasaki2025}.
The boundary hypotheses of that theorem require a separate argument
from the even-periodic estimates proved here.

Finite-volume methods for uniform spectral gaps have a substantial
history, including the criteria of Knabe and Nachtergaele for
frustration-free chains, whose ground states minimize every local
interaction \cite{Knabe1988,Nachtergaele1996}.
Gosset and Mozgunov obtained sharper finite-size thresholds within that
frustration-free setting~\cite{GossetMozgunov2016}.
Recently, Banerjee, Guo and Chowdhury developed semidefinite bounds for
string correlators using an enclosing ground-energy window, including
finite periodic spin-one Heisenberg chains
\cite[Equations~(2), (13) and Figure~4]{BanerjeeGuoChowdhury2026}.
Those observable bounds concern a different quantity from the uniform
spectral gap studied here.
Our finite inputs are thermal traces of small twisted chains and
variational energy estimates from periodic matrix-product vectors, in the
tradition of finitely correlated states \cite{FNW1992}.
The latter use the standard trace-of-matrix-products representation
\cite[Section~2.1]{PerezGarcia2007}; we give the particular integer
matrices and all contractions explicitly.

The finite thermal inputs are rigorous enclosures of matrix-exponential
traces. Validated matrix exponentiation has a long history, including
self-validating Pad\'e methods~\cite{BochevMarkov1989} and interval
polynomial evaluation~\cite{GoldsztejnNeumaier2014}. Our particular
implementation uses a conditioned Poisson polynomial in a contraction.
The factor-two conditioning estimate follows the coefficient argument
of Fox and Glynn~\cite[Proposition~1]{FoxGlynn1988}. Explicit bounds
for the integer evaluation and the final trace enclosure are supplied
in the appendices.

\subsection{How the proof works}

Quantum transfer-matrix methods relate partition functions to spatial
moments~\cite{Suzuki1985}. Our construction uses an ordered bond
expansion of the type developed by Aizenman and
Nachtergaele~\cite[arXiv v1, Section~2.2]{AizenmanNachtergaele1994}, followed by a
direct operator construction on bond histories.

A shifted partition function is a sum of positive physical Boltzmann
weights. We construct a self-adjoint spatial transfer operator whose
$L$th moment is the same partition function. For even $L$, the spatial
moment is also a sum of nonnegative weights. This gives two probability
distributions associated with one partition function. The \emph{purity}
of a distribution $(p_i)$ is $\sum_i p_i^2$; purity close to one means
that a single weight carries almost all of its mass.

Squaring and renormalizing a distribution reduces its purity defect
quadratically once the defect is small. Squaring physical weights
doubles inverse temperature, while squaring spatial weights doubles
length. Exact cancellation identities couple these two improvements.
Starting from two finite purity bounds, both defects then decay doubly
exponentially as length and inverse temperature double together.
Spatial moment inequalities cover every intermediate even length.
Finally, purity of the full physical Gibbs distribution bounds each
excited-to-ground probability ratio and hence the physical energy gap.
Proposition~\ref{prop:bootstrap} states this coupled argument as a general
gap criterion. Its assumptions isolate a reusable criterion
for other Hamiltonian families with the same two partition-function
representations; positivity of the spatial transfer is not required.

There are two finite initializations. One starts the iteration at length
$60$ and inverse temperature $105/4$. The other reaches length $2304$
and inverse temperature $784$ with a smaller purity defect, giving the
stronger eventual bound. They share the transfer construction and the
bootstrap, while using different thermal moments and two lengths of one
matrix-product trial family.
The finite computations involve integer arithmetic, rational interval
bounds and explicit truncation errors; no floating-point eigenvalue
estimate is a premise.

Section~\ref{sec:model} fixes the model and proves its elementary energy
bounds. Section~\ref{sec:transfer} constructs the spatial transfer and
its rotation sectors. Section~\ref{sec:bootstrap} proves the common
purity theorem, and Section~\ref{sec:initialization} establishes the
two finite initializations from explicitly stated thermal and trial
inputs. Section~\ref{sec:assembly} combines these results to prove
Theorem~\ref{thm:main}. Appendices~\ref{sec:thermal120}--\ref{sec:trials}
prove the essential thermal and trial inputs, with exact data and
reproducible finite procedures.
